\documentclass[11pt,letterpaper]{article}
\usepackage[margin=1in]{geometry}
\usepackage{booktabs}
\usepackage{enumitem}
\usepackage{titlesec}
\usepackage{hyperref}
\usepackage{xcolor}
\usepackage{array}
\usepackage{longtable}
\usepackage{tabularx}

\hypersetup{colorlinks=true, linkcolor=blue!60!black, urlcolor=blue!60!black}

\titleformat{\section}{\large\bfseries}{}{0em}{}[\vspace{-0.5em}\rule{\textwidth}{0.4pt}]
\titleformat{\subsection}{\normalsize\bfseries}{}{0em}{}
\titleformat{\subsubsection}{\normalsize\bfseries\itshape}{}{0em}{}

\setlength{\parskip}{6pt}
\setlength{\parindent}{0pt}

\begin{document}

\begin{center}
{\LARGE\bfseries PERSIST Phase I R\&D Plan}\\[6pt]
{\large Topological Diagnostic for Catastrophic Forgetting\\in Production ML Systems}\\[12pt]
{\normalsize\textbf{Axion Deep Labs} \hspace{1.5em} PI: Joshua Gutierrez}\\[2pt]
{\normalsize Target Solicitation: NSF SBIR Phase I}
\end{center}

\vspace{1em}

%--------------------------------------------------------------
\section{Commercial Hypothesis}

Enterprise ML teams retrain and fine-tune models continuously. When models learn new tasks, they catastrophically forget previous capabilities. The current response is heuristic: teams guess which mitigation strategy to apply, tune it by trial and error, and discover failures after deployment. No diagnostic exists that predicts, before committing compute, which strategy will work and how much it will help.

\textbf{The cost of the status quo:} A single mitigation attempt on a 300M-parameter model requires 4--8 GPUs for 8--24 hours of fine-tuning, plus evaluation and debugging. At current cloud H100 pricing (\$3--10/GPU-hour, per published rates from Lambda Labs, RunPod, and major cloud providers), one failed attempt costs \$96--1,920 in compute alone. Teams typically iterate through 3--5 mitigation strategies before finding one that works, spending \$300--9,600 in GPU cost per search. The larger expense is engineering time: a senior ML engineer (\$75--100/hour) spending 2--4 days evaluating each failed attempt adds \$1,200--3,200 per cycle. A typical mitigation search costs \$4,000--15,000 in combined compute and engineering time. For billion-parameter models, these costs scale by an order of magnitude.

\textbf{If topology predicts mitigation benefit at production scale} ($\rho > 0.4$ across 100M+ parameter models and multiple CL methods), \textbf{then:} a single topology scan (estimated cost: \$50--500 depending on model size) replaces the trial-and-error cycle, saving \$4,000--15,000+ per model update. This enables a SaaS diagnostic or SDK integration with immediate ROI for any team managing sequential model updates.

\textbf{Target market segments:}
\begin{itemize}[nosep]
  \item LLM fine-tuning pipelines (every company customizing foundation models)
  \item Medical AI under regulatory update requirements (FDA mandates continuous learning documentation)
  \item Autonomous vehicle perception systems (sequential domain adaptation)
  \item Recommendation systems with evolving product catalogs
\end{itemize}

\textbf{Addressable market:} The MLOps tooling market exceeds \$1B annually. Continual learning failures are a known cost center with no existing diagnostic solution.

%--------------------------------------------------------------
\section{Primary Phase I Objective}

\textbf{Determine whether persistent homology predicts continual learning mitigation benefit at production scale (models > 100M parameters) across multiple mitigation methods.}

\textbf{Central hypothesis:} Basin fragmentation, measured by the H0 persistent homology of loss landscapes, is a structural property that determines the effectiveness of regularization-based continual learning methods, and this relationship persists at production scale.

All Phase I activities test components of this hypothesis. If scale validation and method generalization both succeed, PERSIST has a commercial product. If either fails, we know precisely why and whether a pivot is viable.

%--------------------------------------------------------------
\section{Competitive Landscape}

\textbf{Current industry practice:} Mitigation method selection is heuristic or grid-search based. Teams choose EWC, replay, or architectural methods based on literature or trial and error. No tool predicts which method will work best for a given model before training begins.

\textbf{Existing tools:} CL frameworks (Avalanche, Continuum) implement methods but provide no diagnostic guidance. Model monitoring tools (WhyLabs, Arize) detect drift after deployment but do not predict forgetting before it occurs. Loss landscape visualization (Weights \& Biases) shows surfaces but extracts no predictive features.

\textbf{PERSIST's differentiation:} First predictive diagnostic for continual learning. Topology extracts a structural feature (basin fragmentation) that correlates with mitigation benefit before training begins. Prediction vs.\ detection.

\textbf{Technical novelty:} No prior work has demonstrated a statistically predictive relationship between persistent homology of loss landscapes and mitigation-specific continual learning benefit. Existing loss landscape studies (Li et al., 2018; Keskar et al., 2017) are descriptive: they visualize landscape geometry but do not connect it to actionable outcomes. PERSIST moves from visualization to diagnostic, linking a computable topological invariant (H0) to a decision that practitioners actually face (which mitigation strategy to deploy).

%--------------------------------------------------------------
\section{Preliminary Results (Complete)}

57 configurations (19 architectures $\times$ 3 datasets), all phases complete on models 0.3M--44.7M parameters.

\begin{itemize}[nosep]
  \item \textbf{H0 predicts EWC benefit:} CIFAR-100 $\rho = 0.76$ ($p = 0.0002$), RESISC-45 $\rho = 0.86$ ($p = 2.4\text{e-}6$). Replicates across 2 of 3 datasets.
  \item \textbf{Pooled interaction model:} Dataset moderates the topology-benefit relationship ($p = 0.046$).
  \item \textbf{Conditional forgetting prediction:} On CUB-200, topology rescues prediction where parameter count fails ($p = 0.037$, does not survive Bonferroni).
  \item \textbf{Basin fragmentation hypothesis:} H0 measures landscape fragmentation; regularization benefit scales with fragmentation degree.
\end{itemize}

\textbf{What this establishes:} The signal exists at small scale under controlled conditions. \textbf{What it does not establish:} Whether the signal, the computation, or the commercial value survive at production scale.

%--------------------------------------------------------------
\section{Phase I Technical Plan}

\subsection{Risk 1 (Primary): Scale Survival}

\textbf{Question:} Does the topological signal persist on 100M+ parameter models, or is it a small-model artifact?

All current models are under 45M parameters. Wide networks converge to flatter minima at scale, which could eliminate the H0 variance driving the signal. This is the single question that determines commercial viability.

\textbf{Activities:}
\begin{enumerate}[nosep]
  \item \textbf{Scale ladder} (Months 1--4): Extend WRN width ladder to 100M+, add ViT-Base (86M), ViT-Large (307M). Minimum 5 models above 100M. Compute PH, measure EWC benefit, test H0-benefit correlation at each scale.
  \item \textbf{Signal decay curve} (Months 2--5): Plot correlation strength vs.\ model scale. Identify gradual decay, phase transition, or persistence.
\end{enumerate}

\textbf{Success criteria:} Go if $\rho > 0.4$ ($p < 0.05$) at 100M+. Pivot if $\rho$ 0.2--0.4. No-go if $\rho < 0.2$.

\subsection{Risk 2: Mitigation Method Generalization}

\textbf{Question:} Does topology predict benefit across mechanistically different CL methods?

A commercial diagnostic must work for the method the customer uses. EWC-specific findings have a narrow market.

\textbf{Activities (ordered by cost and information value):}
\begin{enumerate}[nosep]
  \item \textbf{Synaptic Intelligence (SI) replication} (Months 1--3): Run SI on all 57 configs. Cheapest generalization test. SI tracks importance online vs.\ EWC's offline Fisher computation. Positive result generalizes to regularization-based methods broadly.
  \item \textbf{PackNet test} (Months 3--5): Hard weight pruning, mechanistically unrelated to EWC. Positive result means basin fragmentation runs deeper than regularization.
  \item \textbf{Experience replay test} (Months 5--7): Buffer-based, unrelated to landscape geometry. Expected null (supports mechanism if confirmed).
\end{enumerate}

\textbf{Success criteria:} Go if H0 predicts SI + one non-regularization method. Narrow if SI only. No-go if SI null on all datasets.

\textbf{Key decision point (Month 3):} SI results determine whether to invest in scale experiments.

\subsection{Risk 3: Task Sequence Length}

\textbf{Question:} Does topology remain predictive over 10--100+ task sequences?

Production CL involves long curricula. The landscape changes after each task; topology measured once may be irrelevant by task 10.

\textbf{Activities:}
\begin{enumerate}[nosep]
  \item \textbf{5-task and 10-task sequences} (Months 3--7): Split CIFAR-100 into sequential tasks. Compute PH after each boundary. Test prediction horizon.
  \item \textbf{Dynamic topology tracking} (Months 6--8): Test $\Delta$-H0 between tasks as forgetting predictor.
\end{enumerate}

\textbf{Success criteria:} Go if predictive through 5+ tasks. Pivot if 2--3 tasks (still marketable for fine-tuning). No-go if no prediction beyond next task.

\subsection{Risk 4: Computational Tractability}

\textbf{Question:} Can PH extraction cost less than the retraining it prevents?

Ripser scales $O(n^3)$ in simplex count. If extraction at scale exceeds training cost, the tool is impractical regardless of signal quality.

\textbf{Activities:}
\begin{enumerate}[nosep]
  \item \textbf{Scaling benchmark} (Months 1--3): Measure PH time across grid sizes ($50 \times 50$ to $500 \times 500$) and model sizes (1M to 1B). Identify cost crossover.
  \item \textbf{Subsampling fidelity} (Months 2--5): Test whether 5 slices remain sufficient at scale.
  \item \textbf{Distributed PH pipeline} (Months 4--8, if needed): GPU-accelerated Ripser, landmark approximation, distributed cubical complexes.
\end{enumerate}

\textbf{Success criteria:} Go if $<$10\% of training cost at 100M. Viable if 10--50\%. No-go if exceeds training cost with no viable optimization path.

\subsection{Risk 5: Statistical Power and Replication}

\textbf{Question:} Do borderline findings ($p = 0.037$, $p = 0.046$) survive rigorous replication?

\textbf{Activities:}
\begin{enumerate}[nosep]
  \item \textbf{Architecture expansion} (Months 1--6): Add 30+ architectures per dataset. Target 50+ per dataset, 150+ total.
  \item \textbf{Pre-registered replication} (Months 6--8): Register on OSF before expanded analysis. Primary: H0 predicts benefit on 2+ datasets.
\end{enumerate}

\textbf{Success criteria:} Go if $p < 0.01$ pre-registered. Marginal if $p < 0.05$. No-go if $p > 0.05$.

%--------------------------------------------------------------
\section{Customer Discovery (Phase I)}

\textbf{Activities:}
\begin{enumerate}[nosep]
  \item \textbf{20 structured interviews} (Months 1--6): Enterprise ML teams managing sequential model updates. Protocol: current mitigation workflow, pain points, cost of failed updates, willingness to adopt diagnostic tooling. Targets: LLM fine-tuning teams, medical AI companies, AV perception teams, recommendation system engineers.
  \item \textbf{2 letters of interest} (Months 6--9): Contingent on Phase I validation. ``If PERSIST demonstrates X at scale, we would evaluate integration.''
  \item \textbf{Willingness-to-pay threshold}: Hypothesis: \$500--5,000 per model evaluation or \$2,000--10,000/month SaaS.
  \item \textbf{Integration requirements}: Formats (ONNX, PyTorch, TensorFlow), environments (cloud, on-prem, air-gapped), acceptable diagnostic latency.
\end{enumerate}

\textbf{Deliverable:} Customer discovery report with interview summaries, market sizing validation, integration spec, and 2 letters of interest.

%--------------------------------------------------------------
\section{Infrastructure Investment}

\subsection{Equipment: Dedicated Research Workstation (\textasciitilde\$75,000)}

\textbf{Why this equipment is essential to Phase I objectives:} The primary technical objective (Risk 1: Scale Survival) requires training and extracting persistent homology from models exceeding 100M parameters, up to 1B+. These experiments demand sustained access to high-VRAM GPUs with NVLink interconnect for memory pooling across multi-day training runs. While cloud resources are technically viable, they are economically inefficient for the sustained, iterative topology extraction workloads central to this project. Loss landscape computation requires repeated multi-hour GPU sessions with rapid iteration on parameters; cloud scheduling latency, session interruption, and data transfer overhead reduce effective utilization. Owned infrastructure ensures experimental continuity, cost control, and the ability to run experiments around the clock without per-hour cost pressure.

\textbf{Build specification (vendor quotes pending; estimates based on published pricing):}

\begin{center}
\begin{tabular}{@{}p{2.2cm}p{5.5cm}p{2.8cm}p{3cm}@{}}
\toprule
\textbf{Component} & \textbf{Spec} & \textbf{Est.\ Cost} & \textbf{Source} \\
\midrule
GPU & 2$\times$ NVIDIA H100 NVL 94GB (NVLink, 600GB/s) & \$54,000--64,000 & Exxact, PNY \\
CPU & AMD Threadripper PRO 7975WX (32-core, sTR5) & \$3,700--4,000 & Amazon, Micro Center \\
Motherboard & ASUS Pro WS TRX50-SAGE WIFI (PCIe 5.0 x16, CEB) & \$900 & Amazon, Newegg \\
RAM & 512GB DDR5 ECC R-DIMM (8$\times$64GB) & \$1,800 & Vendor pricing \\
Storage & 2TB + 4TB NVMe + 8TB archive & \$600 & Vendor pricing \\
PSU + cooling + case & 1600W Titanium, custom water loop, CEB tower & \$5,500 & Component vendors \\
\midrule
\textbf{Total} & & \textbf{\textasciitilde\$70,000--75,000} & Formal quotes to accompany proposal \\
\bottomrule
\end{tabular}
\end{center}

\textbf{Note:} GPU pricing for the H100 NVL 94GB varies significantly by vendor and procurement channel (published range: \$27,000--\$40,000+ per unit). Final budget figures will reflect formal vendor quotes obtained prior to proposal submission. If pricing exceeds budget, the fallback configuration is 2$\times$ H100 PCIe 80GB (standard, non-NVL variant at lower cost, sacrificing NVLink for PCIe-only inter-GPU communication).

\textbf{Procurement plan:} Vendor quotes from Exxact Corporation, PNY (authorized NVIDIA partner), and BIZON Tech will be solicited during proposal preparation. Equipment will be procured within the first 30 days of the award period. Assembly and validation: 2 weeks. All procurement follows Axion Deep Labs' purchasing policy and NSF cost principles for allowability, allocability, and reasonableness.

\textbf{Usage allocation:} The equipment will be used exclusively for PERSIST objectives during the Phase I performance period except for routine maintenance and system validation. Usage will be logged by experiment ID and linked to the specific technical risk it addresses. A pre-owned RTX 4090 workstation (not a Phase I cost) handles ancillary small-model experiments, ensuring the H100 system is reserved for scale-dependent Phase I objectives.

\textbf{Cost comparison vs.\ cloud:} Published H100 cloud rental rates range from \$3.00--\$10.60/GPU-hour depending on provider, commitment terms, and availability (sources: Lambda Labs, RunPod, Jarvislabs 2026 pricing guides). At a conservative midpoint of \$5/GPU-hour for two GPUs (\$10/hr), the 5,000 GPU-hour Phase I experiment set would cost \$25,000--50,000 in cloud rental with no residual asset. The dedicated workstation reaches cost parity during Phase I at sustained utilization and provides zero-marginal-cost compute for Phase II development.

\subsection{Compute Prioritization}

All experiments run on owned hardware. No cloud dependency, no allocation uncertainty. Prioritization reflects experimental logic, not budget constraints:

\begin{center}
\begin{tabular}{@{}clcl@{}}
\toprule
\textbf{Priority} & \textbf{Experiment} & \textbf{Est.\ GPU-hrs} & \textbf{Decision it enables} \\
\midrule
1 & SI replication (57 configs) & \textasciitilde200 & Method generalization go/no-go (Mo.\ 3) \\
2 & Scale ladder to 100M (5 models) & \textasciitilde500 & Signal at production scale \\
3 & Architecture expansion (+30) & \textasciitilde600 & Statistical power for replication \\
4 & 5-task sequences (top 10 archs) & \textasciitilde300 & Task sequence viability \\
5 & PH scaling benchmark & \textasciitilde100 & Tractability assessment \\
6 & PackNet + replay (57 configs each) & \textasciitilde400 & Cross-method generalization \\
7 & ViT-Large (307M) + 1B model & \textasciitilde1,500 & Deep scale validation \\
8 & 10-task sequences + dynamic topology & \textasciitilde500 & Long-sequence prediction \\
9 & Pre-registered replication (150+ configs) & \textasciitilde900 & Confirmatory analysis \\
\midrule
\textbf{Total} & & \textbf{\textasciitilde5,000} & Full Phase I experiment set \\
\bottomrule
\end{tabular}
\end{center}

Priorities 1--5 (minimum viable set, \textasciitilde1,700 hours) answer the core questions. Priorities 6--9 deepen and confirm. With dedicated hardware running continuously, the full 5,000-hour set completes comfortably within 9 months (\textasciitilde19 hours/day average utilization).

\textbf{If a go/no-go gate fails early} (e.g., SI null at Month 3), remaining Phase I effort redirects to investigating the failure mechanism (publishable, advancing the field) and documenting negative results for the community.

%--------------------------------------------------------------
\section{Timeline and Milestones}

\begin{center}
\begin{tabular}{@{}clll@{}}
\toprule
\textbf{Month} & \textbf{Activity} & \textbf{Deliverable} & \textbf{Gate} \\
\midrule
1--2 & SI replication + PH benchmark + interviews & SI data, cost model & \\
3 & \textbf{SI go/no-go} & SI results report & \textbf{No-go if null on all datasets} \\
2--4 & Scale ladder + architecture expansion & Signal at scale & \\
5 & \textbf{Scale go/no-go} & Signal decay curve & \textbf{No-go if $\rho < 0.2$ at 100M} \\
3--5 & PackNet + 5-task sequences & Method + sequence data & \\
6 & Pre-register on OSF & Registration & \\
6--8 & Pre-registered analysis + ViT-Large & Confirmatory results & \\
8 & \textbf{Replication go/no-go} & Pre-registered analysis & \textbf{No-go if $p > 0.05$} \\
7--9 & Dynamic topology + distributed PH + LOIs & Prototype pipeline & \\
9 & Final analysis + Phase I report & All deliverables & Phase II decision \\
\bottomrule
\end{tabular}
\end{center}

%--------------------------------------------------------------
\section{Budget Overview}

\begin{center}
\begin{tabular}{@{}lrr@{}}
\toprule
\textbf{Category} & \textbf{Est.\ Cost} & \textbf{\% of Phase I} \\
\midrule
Equipment (research workstation) & \$75,000 & 27\% \\
Personnel (PI full-time + Senior Personnel part-time) & \$118,500 & 43\% \\
NSF I-Corps (commercialization training) & \$25,000 & 9\% \\
TABA (Technical and Business Assistance) & \$6,500 & 2\% \\
Customer discovery (travel, interviews) & \$15,000 & 6\% \\
Materials, supplies, indirect & \$35,000 & 13\% \\
\midrule
\textbf{Total} & \textbf{\textasciitilde\$275,000} & \textbf{100\%} \\
\bottomrule
\end{tabular}
\end{center}

\textbf{I-Corps:} Dedicated commercialization training for the PI, including customer discovery methodology, market validation frameworks, and mentor engagement through the NSF I-Corps program. The PI commits to completing the NSF I-Corps customer discovery curriculum during Phase I. Directly supports the 20-interview customer discovery plan and go-to-market readiness.

\textbf{TABA:} Third-party technical and business advisory services for IP strategy assessment, market sizing validation, and Phase II commercialization planning. Supports the open-core IP model and pricing strategy development.

\textbf{Equipment justification:} See Infrastructure Investment section. The dedicated workstation is essential to achieving Phase I technical objectives that cannot be met with cloud resources at this budget level. Cloud rental for the required 5,000 GPU-hours would consume \$30,000--60,000 with no residual asset. The workstation is allocated primarily (\textasciitilde90\%) to this award's experiments during Phase I, with usage logged by experiment for audit purposes.

%--------------------------------------------------------------
\section{Personnel and Team Capacity}

\begin{center}
\begin{tabular}{@{}lll@{}}
\toprule
\textbf{Role} & \textbf{Effort} & \textbf{Responsibility} \\
\midrule
PI (Joshua Gutierrez) & Full-time & Experimental design, PH computation, statistical analysis \\
Senior Personnel (Crystal Gutierrez) & Part-time & ML pipeline development, CL method implementation \\
\bottomrule
\end{tabular}
\end{center}

The PI has already executed all 57 configurations (19 architectures $\times$ 3 datasets, 7 phases each) end-to-end without external support, including custom dataset pipelines, PH extraction, statistical analysis, and the pooled interaction model. Senior Personnel independently built and evaluated ML models on industrial datasets at Bayer through the Purdue Data Mine and currently teaches data analysis methods as an adjunct professor at NMSU. This team has demonstrated capacity to manage large experimental pipelines with minimal overhead.

%--------------------------------------------------------------
\section{Intellectual Property Strategy}

\textbf{Open core model:}
\begin{itemize}[nosep]
  \item \textbf{Open:} Research methodology, preliminary benchmark data, and core PH extraction code. Builds scientific credibility and drives adoption.
  \item \textbf{Proprietary:} Scalable distributed PH implementation (Phase I R\&D output), production-scale benchmark dataset during commercialization window, and method-selection algorithms.
  \item \textbf{Patent potential:} Novel distributed PH extraction methods developed under Risk 4, if applicable. Provisional patent filings during Phase I.
\end{itemize}

This resolves the open-science vs.\ commercial-defensibility tension: the science is reproducible, but the production engineering is proprietary.

%--------------------------------------------------------------
\section{Phase I Deliverables}

\begin{center}
\begin{tabular}{@{}p{4.5cm}p{10cm}@{}}
\toprule
\textbf{Deliverable} & \textbf{Description} \\
\midrule
Open benchmark dataset & Topology metrics, forgetting curves, and mitigation benefit scores across 150+ configurations \\
Scalable PH prototype & Python package for topology extraction, tested on 300M+ parameter models \\
Pre-registered replication & OSF-registered confirmatory analysis (or documented null) \\
Peer-reviewed submission & NeurIPS/ICML target reporting Phase I results \\
ArXiv preprint & Preliminary results (in progress, pending endorsement) \\
Customer discovery report & 20 interviews, integration spec, WTP data, market sizing \\
2 letters of interest & Enterprise ML teams, contingent on validation \\
Phase II plan & Product development and go-to-market (if Phase I succeeds) \\
\bottomrule
\end{tabular}
\end{center}

%--------------------------------------------------------------
\section{Broader Impacts: Trustworthy AI}

Catastrophic forgetting is a reliability and safety problem:

\begin{itemize}[nosep]
  \item \textbf{Defense systems} learning sequentially must not forget previous threats when adapting to new ones. A predictive diagnostic prevents silent capability degradation.
  \item \textbf{Medical AI} under FDA requirements must demonstrate that updates do not compromise existing accuracy. Topology provides a pre-deployment safety check.
  \item \textbf{Critical infrastructure AI} under continuous distribution shift needs guaranteed retention.
\end{itemize}

Phase II will explore regulatory documentation support in partnership with domain experts. No clinical or defense deployment is proposed during Phase I.

%--------------------------------------------------------------
\section{Phase II Vision (Contingent on Phase I)}

\begin{itemize}[nosep]
  \item \textbf{Product:} SaaS diagnostic and SDK. Input: pretrained model + task description. Output: recommended mitigation strategy, expected benefit, confidence score.
  \item \textbf{Integration:} Plugins for CL frameworks (Avalanche, Continuum), MLOps platforms (MLflow, W\&B), cloud pipelines (SageMaker, Vertex AI).
  \item \textbf{Revenue:} Per-evaluation (\$500--5,000) or monthly SaaS (\$2,000--10,000/month).
  \item \textbf{Infrastructure continuity:} Phase I workstation serves as the development and pilot deployment platform for Phase II, eliminating additional infrastructure costs during product engineering.
  \item \textbf{Activities:} Product engineering, beta with LOI partners, regulatory documentation framework, go-to-market.
\end{itemize}

%--------------------------------------------------------------
\section{Falsification Commitment}

\begin{center}
\begin{tabular}{@{}p{4cm}cp{7.5cm}@{}}
\toprule
\textbf{Decision point} & \textbf{Timeline} & \textbf{Consequence} \\
\midrule
SI null on all datasets & Month 3 & Finding is EWC-specific. Commercialization path will not proceed to Phase II absent positive validation. \\
Signal vanishes at 100M & Month 5 & Small-model artifact. Commercialization path will not proceed to Phase II absent positive validation. \\
Pre-registered replication fails & Month 8 & Preliminary results not confirmed. Commercialization path will not proceed to Phase II absent positive validation. \\
PH intractable, no workaround & Month 9 & Pivot to approximate topology methods or commercialization path will not proceed to Phase II. \\
\bottomrule
\end{tabular}
\end{center}

Commercialization path will not proceed to Phase II absent positive validation. Phase I funds understanding why the signal fails (publishable, advancing the field) but product development halts without confirmed technical outcomes.

\end{document}
